\oldpage[IV-3]{173}
\section{Estudo das fibras dos morfismos planos de apresentação finita}
\label{section:IV.12}

\setcounter{subsection}{-1}
\subsection{Introdução}
\label{subsection:IV.12.0}

\emph{Utilizaremos em toda esta seção as notações gerais de}
\sref{IV.9.4.1}.

\begin{env}[12.0.1]
\label{IV.12.0.1}
Dado um morfismo localmente de apresentação finita $f:X\to Y$, vimos
\sref{IV.9.9} que, para certas propriedades locais $\mathbf P$ de
pré-esquemas sobre corpos ou de Módulos sobre esses pré-esquemas, o
conjunto $E$ dos $x\in X$ tais que a propriedade $\mathbf P$ se verifica,
para a fibra $X_{f(x)}$, no ponto $x$ dessa fibra, é localmente
construtível em $X$. Propomo-nos mostrar que, para a maioria dessas
propriedades, se supusermos além disso que o morfismo $f$ é plano,
então o conjunto $E$ é até aberto em $X$. Do mesmo modo, em
\sref{IV.9.2}–\sref{IV.9.8}, mostramos que, se $f$ é de apresentação
finita e se, desta vez, $\mathbf P$ designa certas propriedades globais
de pré-esquemas sobre corpos ou de Módulos sobre esses pré-esquemas, o
conjunto $F$ dos $y\in Y$ tais que a propriedade $\mathbf P$ se
verifica para a fibra $X_y$ é localmente construtível em $Y$.
Mostraremos que, se supusermos além disso que o morfismo $f$ é
próprio e plano, $F$ é até aberto em $Y$.
\end{env}

\begin{env}[12.0.2]
\label{IV.12.0.2}
O método geral de demonstração das propriedades em questão compreende
três etapas. Primeiro, reduz-se ao caso em que $Y$ é afim e $X$ é de
apresentação finita; depois:
\begin{enumerate}[label=\emph{\Alph*)}]
\item Com o auxílio de \sref{IV.8.9.1} (e, eventualmente, de outros
resultados da \S~8) e de \sref{IV.11.2.6}, reduz-se ao caso em que
$X$ e $Y$ são noetherianos.
\item Aplicam-se os resultados da \S~9 recordados em
\sref{IV.12.0.1}, que demonstram que $E$ (resp.\ $F$) é construtível.
\item Para ver que $E$ é aberto, basta, em virtude de
\sref[0]{0.9.2.5}, demonstrar que, se $x\in E$, toda generização
$x'$ de $x$ também pertence a $E$. Utilizando
\sref[II]{II.7.1.7}, vemos que, como $Y$ é noetheriano, existem o
espectro $Y_1$ de um anel de valoração discreta e um morfismo
$h:Y_1\to X$ tais que, se $y_1$ (resp.\ $y'_1$) é o ponto fechado
(resp.\ o ponto genérico) de $Y_1$, temos $h(y_1)=x$ e
$h(y'_1)=x'$. Faz-se então a mudança de base
$g=f\circ h:Y_1\to Y$; levando em conta os resultados das
\S\S~4 e 6 sobre os pré-esquemas localmente noetherianos sobre
corpos e as mudanças do corpo de base, reduz-se a demonstrar a
afirmação em questão para um ponto $x_1$ de
$X_1=X\times_Y Y_1$ situado sobre $y_1$ e uma generização
$x'_1$ de $x_1$ situada sobre $y'_1$. Como
$Y_1=\operatorname{Spec}(A)$, onde $A$ é um anel de valoração
discreta, com uniformizante $t$, deve-se finalmente demonstrar,
para um $A$-módulo $M$, uma propriedade de $M$, sabendo que
$M/tM$ possui a mesma propriedade e que $t$ é $M$-regular
(o que decorre da hipótese de planicidade); para isso,
utilizam-se os resultados de \sref{IV.3.4} e
\sref{IV.5.12}. Procede-se do mesmo modo para o conjunto
$F\subset Y$.
\end{enumerate}
\end{env}

\oldpage[IV-3]{174}

\subsection{Propriedades locais das fibras de um morfismo plano localmente de apresentação finita}
\label{subsection:IV.12.1}

\begin{theorem}[12.1.1]
\label{IV.12.1.1}
Sejam $f:X\to Y$ um morfismo localmente de apresentação finita,
$\sh{F}$ um $\sh{O}_X$-Módulo $f$-plano e de apresentação finita,
$\Phi$ uma parte finita de $\mathbf Z\cup\{\pm\infty\}$ e $k$ um
inteiro. As seguintes partes de $X$ são abertas:
\begin{enumerate}[label=(\roman*)]
\item O conjunto dos $x\in X$ tais que as dimensões dos ciclos
primos associados a $\sh{F}_{f(x)}$ e que contêm $x$ sejam
elementos de $\Phi$.
\item O conjunto dos $x\in X$ tais que todos os ciclos primos
associados a $\sh{F}_{f(x)}$ e que contêm $x$ tenham dimensão $k$.
\item O conjunto dos $x\in X$ tais que $x$ não pertença a nenhum
ciclo primo imerso associado a $\sh{F}_{f(x)}$.
\item O conjunto dos $x\in X$ tais que $\sh{F}_{f(x)}$ seja
equidimensional no ponto $x$ e possua a propriedade $(S_k)$ nesse
ponto \sref{IV.5.7.2}.
\item O conjunto dos $x\in X$ tais que
$\operatorname{coprof}((\sh{F}_{f(x)})_x)\leq k$
\sref[0]{0.16.4.9}.
\item O conjunto dos $x\in X$ tais que $\sh{F}_{f(x)}$ seja um
$\sh{O}_{X_{f(x)}}$-Módulo de Cohen--Macaulay no ponto $x$
\sref{IV.5.7.1}.
\item O conjunto dos $x\in X$ tais que $\sh{F}_{f(x)}$ seja
geometricamente reduzido no ponto $x$ \sref{IV.4.6.22}.
\item O conjunto dos $x\in X$ tais que $\sh{F}_{f(x)}$ seja
geometricamente pontualmente integral no ponto $x$
\sref{IV.4.6.22}.
\end{enumerate}
\end{theorem}

\begin{proof}
Como as questões são locais em $X$, reduz-se primeiro ao caso em que
$X=\operatorname{Spec}(B)$ e $Y=\operatorname{Spec}(A)$ são afins e
$f$ é um morfismo de apresentação finita. Existem então um
subanel noetheriano $A_0$ de $A$, um $A_0$-pré-esquema de tipo
finito $X_0$ e um $\sh{O}_{X_0}$-Módulo coerente $\sh{F}_0$ tais
que $\sh{F}_0\otimes_{A_0}A$ é isomorfo a $\sh{F}$
\sref{IV.8.9.1}; além disso, em virtude de \sref{IV.11.2.6},
pode-se supor que $\sh{F}_0$ é plano sobre
$Y_0=\operatorname{Spec}(A_0)$. Se $h:X\to X_0$ é a projeção
canônica, o conjunto $E$ dos pontos $x\in X$ nos quais se verifica
uma das propriedades (i) a (viii) é igual a $h^{-1}(E_0)$, onde
$E_0$ é o conjunto dos $x_0\in X_0$ nos quais se verifica a
propriedade correspondente para $\sh{F}_0$. Isso resulta, para
as propriedades (i) a (iii), de \sref{IV.4.2.7}; para (iv) a
(vi), de \sref{IV.6.7.1}; e para (vii) e (viii), de
\sref{IV.4.7.11}.

Assim, reduz-se ao caso em que $A$ e $B$ são noetherianos, $f$ é
de tipo finito e $\sh{F}=\widetilde M$, sendo $M$ um $B$-módulo
de tipo finito. Por \sref{IV.9.9.2}, sabe-se que $E$ é construtível
para todas as propriedades consideradas; resta, em cada caso, a
etapa~C) de \sref{IV.12.0.2}, na qual se deve demonstrar que $E$
é estável por generização.

\begin{env}[12.1.1.1]
\label{IV.12.1.1.1}
Comecemos pelo caso~(iii), o mais simples. Seja $x'\ne x$ uma
generização de $x$ em $X$. Ponhamos $y=f(x)$ e $y'=f(x')$.
Existem o espectro $Y_1$ de um anel de valoração discreta e um
morfismo $u:Y_1\to X$ tais que, se $y_1$ e $y'_1$ são,
respectivamente, o ponto fechado e o ponto genérico de $Y_1$,
temos $u(y_1)=x$ e $u(y'_1)=x'$ \sref[II]{II.7.1.7}.
Ponhamos $g=f\circ u:Y_1\to Y$. Se
$X_1=X\times_Y Y_1$, existe uma $Y_1$-seção
$u_1:Y_1\to X_1$ tal que $u=g_1\circ u_1$, onde
$g_1:X_1\to X$ é a projeção canônica. Se
$x_1=u_1(y_1)$ e $x'_1=u_1(y'_1)$, temos
$g_1(x_1)=x$ e $g_1(x'_1)=x'$, e $x'_1$ é uma
generização de $x_1$. Aplicando novamente
\sref{IV.4.2.7}, vê-se
\oldpage[IV-3]{175}
que se reduz ao caso em que $Y=\operatorname{Spec}(A)$ é o
espectro de um anel de valoração discreta, sendo $y$ o ponto
fechado e $y'$ o ponto genérico de $Y$. Se $t$ é um
uniformizante de $A$, a hipótese de que $\sh{F}$ é $f$-plano
implica que $t$ é $\sh{F}$-regular \sref[0]{0.6.3.4}.
Por hipótese, nenhum dos ciclos primos imersos associados a
$\sh{F}/t\sh{F}$ contém $x$. Então, de
\sref{IV.3.4.4}, resulta que $x$ não pertence a nenhum ciclo
primo imerso associado a $\sh{F}$, e o mesmo vale para toda
generização de $x$. Em particular, $x'$ não pertence a nenhum
ciclo primo imerso associado a $\sh{F}$, nem
\emph{a fortiori} a nenhum dos associados a $\sh{F}_{y'}$
(pois $X_{y'}$ é um subpré-esquema induzido em um aberto de $X$).
\end{env}

\begin{env}[12.1.1.2]
\label{IV.12.1.1.2}
Consideremos agora os casos~(iv) e~(v) (o caso~(vi) não passa de
um caso particular de~(v)). Procede-se como acima, usando
\sref{IV.6.7.1} em lugar de \sref{IV.4.2.7}, e reduz-se ao caso
em que $Y$ é o espectro de um anel de valoração discreta $A$.

O anel $C=\sh{O}_{X,x}$ é então um localizado de uma
$A$-álgebra de tipo finito; portanto, é catenário
\sref{IV.5.6.4}. Por hipótese, o $C$-módulo $M_x/tM_x$
satisfaz $(S_k)$ e é equidimensional (resp.\ tem-se
$\operatorname{coprof}(M_x/tM_x)\leq k$). Deduz-se, pois, de
\sref{IV.5.12.2} (resp.\ de \sref[0]{0.16.4.10}) que $M_x$
satisfaz $(S_k)$ e é equidimensional (resp.\ que
$\operatorname{coprof}(M_x)\leq k$), pois $t$ é $M_x$-regular
e pertence ao ideal maximal de $C$. Daí as conclusões, já que
$x'$ é uma generização de $x$ e
$(\sh{F}_{y'})_{x'}=\sh{F}_{x'}$.
\end{env}

\begin{env}[12.1.1.3]
\label{IV.12.1.1.3}
Passemos aos casos~(vii) e~(viii). Podemos substituir $Y$ pelo
subpré-esquema integral cujo espaço subjacente é
$\overline{\{y'\}}$ e $X$ por
$f^{-1}(\overline{\{y'\}})$ sem alterar as fibras em $y$ e
$y'$; podemos, portanto, supor que $A$ é integral, com corpo
de frações $K=k(y')$. Sabe-se por \sref{IV.4.5.11} e
\sref{IV.4.7.8} que existe uma extensão finita $K'$ de $K$
tal que, para o $(\sh{O}_X\otimes_A K')$-Módulo
$\sh{F}'=\sh{F}\otimes_A K'$, os ciclos primos associados
são geometricamente irredutíveis e os comprimentos geométricos
de $\sh{F}'$ nos pontos maximais $z'_i$ de seu suporte são,
respectivamente, iguais aos comprimentos de
$\sh{F}'_{z'_i}$ nesses pontos. Assim, levando em conta
\sref{IV.4.2.7}, equivale dizer que $\sh{F}$ é
geometricamente reduzido (resp.\ geometricamente
pontualmente integral) em um ponto $x'$ de
$X_{y'}=X\otimes_A K$, ou dizer que $\sh{F}'$ é
\emph{reduzido} (resp.\ \emph{integral}) nos pontos de
$X\otimes_A K'$ situados sobre $x'$.

Seja $A'$ uma $A$-álgebra finita cujo corpo de frações é
$K'$. Por \sref{IV.4.7.11}, em todo ponto de
$X\otimes_A A'$ situado sobre $x$, o Módulo
$\sh{F}\otimes_A A'$ possui a propriedade~(vii)
(resp.~(viii)). Substituindo $A$ por $A'$ e $X$ por
$X\otimes_A A'$, vê-se que basta demonstrar que $\sh{F}$
é \emph{reduzido} (resp.\ \emph{integral}) em toda generização
$x'$ de $x$. Procede-se então como em
\sref{IV.12.1.1.1}, utilizando desta vez
\sref{IV.4.7.11}, e reduz-se novamente ao caso em que $A$
é um anel de valoração discreta, $y$ é o ponto fechado,
$y'$ o ponto genérico de $Y$ e o uniformizante $t$ de $A$
é um elemento $M$-regular. Como, por hipótese,
$\sh{F}/t\sh{F}$ é reduzido (resp.\ integral) no ponto $x$,
resulta de \sref{IV.3.4.6} (resp.\ de
\sref{IV.3.4.5}) que $\sh{F}$ é reduzido (resp.\ integral)
em $x$ e, portanto, também em uma vizinhança de $x$,
em particular em $x'$. Isso conclui a demonstração nos
casos~(vii) e~(viii).
\end{env}

\begin{env}[12.1.1.4]
\label{IV.12.1.1.4}
Resta examinar os casos~(i) e~(ii). Substituindo $Y$ pelo
subpré-esquema integral cujo espaço subjacente é
$\overline{\{y'\}}$, podemos restringir-nos ao caso em que $Y$
é irredutível e $y'$ é seu ponto genérico. Seja $z'$ um ponto
genérico de um ciclo primo associado a $\sh{F}_{y'}$
\oldpage[IV-3]{176}
que contém $x'$, e seja $Z'$ o fecho de $z'$ em $X$, de modo
que $x\in Z'$. Para tratar o caso~(i), basta demonstrar a

\begin{proposition}[12.1.1.5]
\label{IV.12.1.1.5}
Sejam $Y$ um pré-esquema localmente noetheriano,
$f:X\to Y$ um morfismo localmente de tipo finito e
$\sh{F}$ um $\sh{O}_X$-Módulo coerente e $f$-plano;
para todo $y\in Y$, ponhamos
$\sh{F}_y=\sh{F}\otimes_{\sh{O}_Y}k(y)$.
Seja $z'$ um ponto de $X$, seja $y'=f(z')$ e
suponhamos que $z'\in\operatorname{Ass}(\sh{F}_{y'})$.
Sejam $Z'$ (resp.\ $Y'$) o subpré-esquema reduzido
de $X$ (resp.\ de $Y$) cujo espaço subjacente é o fecho
de $\{z'\}$ em $X$ (resp.\ de $\{y'\}$ em $Y$).
Então, para todo $y\in f(Z')$, as dimensões de todas
as componentes irredutíveis de $Z'_y$ são iguais a
$\dim(Z'_{y'})$ (dimensão do fecho de $z'$ em
$X_{y'}$) (o que exprimiremos mais precisamente em
\sref{IV.13.2.2}, dizendo que a restrição
$Z'\to Y'$ de $f$ é um morfismo equidimensional);
além disso, em todo ponto maximal $z$ de $Z'_y$,
tem-se $z\in\operatorname{Ass}(\sh{F}_y)$.
\end{proposition}

Para isso, reduziremos ao caso em que $Y$ é o espectro
de um anel de valoração discreta. Tomemos o espectro
$Y_1$ de um anel de valoração discreta e um morfismo
$h:Y_1\to X$ tais que $h(y_1)=z$ e $h(y'_1)=z'$,
sendo $y_1$ e $y'_1$, respectivamente, o ponto fechado
e o ponto genérico de $Y_1$ \sref[II]{II.7.1.7}.
Sejam $g=f\circ h:Y_1\to Y$,
$X_1=X\times_Y Y_1$, e sejam
$f_1:X_1\to Y_1$ e $g_1:X_1\to X$ as projeções
canônicas. Existe uma $Y_1$-seção $h_1:Y_1\to X_1$
tal que $h=g_1\circ h_1$ \sref[I]{I.3.3.14}.
Se $Z'_1=Z'_{y'}\otimes_{k(y')}k(y'_1)$,
sabe-se por \sref{IV.4.2.6} que as componentes
irredutíveis de $Z'_1$ dominam $Z'_{y'}$.
Seja $z'_1$ o ponto genérico de uma dessas componentes
que contém $h_1(y'_1)$, de modo que
$f_1(z'_1)=y'_1$ e $g_1(z'_1)=z'$.
Como $h_1(y_1)$ é uma especialização de
$h_1(y'_1)$, é \emph{a fortiori} uma especialização
de $z'_1$. Seja $z_1$ um ponto genérico de uma das
componentes irredutíveis de
$\overline{\{z'_1\}}\cap(X_1)_{y_1}$ que contêm
$h_1(y_1)$. Então $g_1(z_1)$ é uma especialização
de $g_1(z'_1)=z'$ em $X$, e
$z=h(y_1)=g_1(h_1(y_1))$ é uma especialização de
$g_1(z_1)$. Mas, como $f(g_1(z_1))=y$ e $z$ é
um ponto genérico de $Z'_y$, tem-se necessariamente
$z=g_1(z_1)$.

Suponhamos agora demonstrado que, se se põe
$\sh{F}_1=\sh{F}\otimes_{\sh{O}_Y}\sh{O}_{Y_1}$,
então $z_1\in\operatorname{Ass}((\sh{F}_1)_{y_1})$.
Resultará de \sref{IV.4.2.7} que
$z\in\operatorname{Ass}(\sh{F}_y)$; além disso,
as dimensões de
$\overline{\{z_1\}}\cap(X_1)_{y_1}$ e de
$\overline{\{z\}}\cap X_y$ são iguais, e também
são iguais as dimensões de
$\overline{\{z'_1\}}\cap(X_1)_{y'_1}$ e de
$\overline{\{z'\}}\cap X_{y'}$
\sref{IV.4.2.7}. Reduzimo-nos, portanto, como
anunciado, ao caso em que $Y$ é o espectro de um
anel de valoração discreta $A$, sendo $y$ e $y'$,
respectivamente, seu ponto fechado e seu ponto
genérico.

Nesse caso, como $z'\in\operatorname{Ass}(\sh{F}_{y'})$,
tem-se \emph{a fortiori}
$z'\in\operatorname{Ass}(\sh{F})$. Se $t$ é um
uniformizante de $A$, então $t$ é $\sh{F}$-regular
pela planicidade; logo, por \sref{IV.3.4.3},
\[
z\in\operatorname{Ass}(\sh{F}/t\sh{F})
=\operatorname{Ass}(\sh{F}_y).
\]
Quanto à afirmação relativa às dimensões, ela resulta
de \sref{IV.7.1.13}, aplicado a um subpré-esquema
fechado de $X$ cujo espaço subjacente é $Z'$.

Abordemos, por fim, o caso~(ii) de \sref{IV.12.1.1}.
Com as mesmas notações, é preciso demonstrar que, se
todos os ciclos primos associados a $\sh{F}_y$ e que
contêm $x$ têm dimensão $k$, o mesmo vale para todos
os ciclos primos associados a $\sh{F}_{y'}$ e que
contêm $x'$. Ora, acabamos de ver que \emph{todo}
ciclo primo associado a $\sh{F}_{y'}$ e que contém
$x'$ tem a mesma dimensão que \emph{um} dos ciclos
primos associados a $\sh{F}_y$ e que contêm $x$.
Isso demonstra~(ii).
\end{env}
\end{proof}

\begin{remarks}[12.1.2]
\label{IV.12.1.2}
\begin{enumerate}[label=(\roman*)]
\item Nas condições de \sref{IV.12.1.1.5}, com
$\sh{F}=\sh{O}_X$ (de modo que $f$ é plano), em geral
não se pode afirmar que a restrição $Z'\to Y'$ de $f$
seja um morfismo plano, nem sequer um morfismo aberto.
Isso é mostrado pelo exemplo \sref{IV.6.5.5},
\oldpage[IV-3]{177}
(ii), no qual se toma como $Z'$ uma das duas componentes
irredutíveis de $X=\operatorname{Spec}(B)$. É imediato
que esse morfismo restrição não é aberto nos pontos
de $Z'$ situados sobre o
\textquotedblleft{}ponto duplo\textquotedblright{}
de $Y'=Y$.
\item Sob as hipóteses de \sref{IV.12.1.1.5}, com
$\sh{F}=\sh{O}_X$, não se pode enfraquecer a condição
\textquotedblleft{}$f$ é plano\textquotedblright{}
substituindo-a por
\textquotedblleft{}$f$ é universalmente aberto\textquotedblright{}
(cf.\ \sref{IV.2.4.6}), como veremos adiante
em um exemplo \sref{IV.14.4.10},~(i).
\end{enumerate}
\end{remarks}

\begin{corollary}[12.1.3]
\label{IV.12.1.3}
Sob as hipóteses de \sref{IV.12.1.1}, a função
$x\mapsto\operatorname{coprof}((\sh{F}_{f(x)})_x)$
é semicontínua superiormente em $X$, e a função
$x\mapsto\operatorname{prof}_x^*(\sh{F}_{f(x)})$
\sref{IV.10.8.1} é semicontínua inferiormente em $X$.
\end{corollary}

\begin{proof}
A primeira afirmação é apenas uma formulação equivalente
de \sref{IV.12.1.1},~(v). Para demonstrar a segunda,
podemos primeiro, levando em conta \sref{IV.10.8.7} e
\sref{IV.10.8.8}, reduzir-nos, como em
\sref{IV.12.1.1}, ao caso em que $Y$ é o espectro
de um anel de valoração discreta $A$, com ponto
fechado $y$ e ponto genérico $y'$, sendo
$X=\operatorname{Spec}(B)$ afim e
$\sh{F}=\widetilde M$, onde $M$ é um $B$-módulo
de tipo finito. Como $\sh{F}$ é $f$-plano, toda
componente irredutível de $\operatorname{Supp}(\sh{F})$
que contém $x\in X_y$ domina $Y$
\sref{IV.2.3.4}; em outras palavras, seu ponto
genérico $z'$ pertence a $X_{y'}$. As componentes
irredutíveis $Z$ de $\operatorname{Supp}(\sh{F})$
que contêm uma generização $x'$ de $x$ pertencente
a $X_{y'}$ são, portanto, exatamente as que contêm
$x$; além disso, por \sref{IV.12.1.1.5}, tem-se
$\dim(Z_{y'})=\dim(Z_y)$, donde, se
$T=\operatorname{Supp}(\sh{F})$,
$\dim_x(T_y)=\dim_{x'}(T_{y'})$.

Ademais, para um uniformizante $t$ de $A$, vimos em
\sref{IV.12.1.1},~(v) que
$\operatorname{coprof}(M_x)=
\operatorname{coprof}(M_x/tM_x)$; como, por outro lado,
$\operatorname{coprof}(M_{x'})\leq
\operatorname{coprof}(M_x)$ \sref{IV.6.11.5},
resulta que
\[
\operatorname{coprof}((\sh{F}_{y'})_{x'})
\leq\operatorname{coprof}((\sh{F}_y)_x).
\]
A relação
$\operatorname{prof}_x^*(\sh{F}_y)\leq
\operatorname{prof}_{x'}^*(\sh{F}_{y'})$
decorre então de \sref{IV.10.8.7}.
\end{proof}

\begin{remark}[12.1.4]
\label{IV.12.1.4}
Deduz-se também de \sref{IV.12.1.1},~(i) que a função
$x\mapsto\dim_x(X_{f(x)})$ é semicontínua
superiormente em $X$, mas veremos mais adiante
\sref{IV.13.1.3} que essa propriedade vale mesmo sem
supor que $f$ seja plano.
\end{remark}

\begin{corollary}[12.1.5]
\label{IV.12.1.5}
Sejam $Y$ um pré-esquema localmente noetheriano,
$f:X\to Y$ um morfismo localmente de tipo finito,
$\sh{F}$ um $\sh{O}_X$-Módulo coerente e $f$-plano,
$\sh{G}$ um $\sh{O}_Y$-Módulo coerente, $x$ um ponto
de $X$ e $y=f(x)$. Suponhamos que $\sh{G}$ possua
a propriedade $(S_k)$ no ponto $y$ e que
$\sh{F}_y$ possua a propriedade $(S_k)$ no ponto
$x$ e seja equidimensional nesse ponto. Então:
\begin{enumerate}[label=(\roman*)]
\item $\sh{F}\otimes_{\sh{O}_Y}\sh{G}$ possui a
propriedade $(S_k)$ no ponto $x$.
\item Se, além disso, $Y$ é localmente imersível em
um esquema regular ou é um pré-esquema excelente
\sref{IV.7.8.5}, existe uma vizinhança de $x$ em
$X$ na qual
$\sh{F}\otimes_{\sh{O}_Y}\sh{G}$ possui a
propriedade $(S_k)$.
\end{enumerate}
\end{corollary}

\begin{proof}
Com efeito, por \sref{IV.12.1.1},~(iv), existe uma
vizinhança aberta $V$ de $x$ tal que, para todo
$x'\in V$, $\sh{F}_{f(x')}$ satisfaz $(S_k)$ no
ponto $x'$. Por outro lado, $\sh{G}$ satisfaz
$(S_k)$ em todos os pontos de
$Y'=\operatorname{Spec}(\sh{O}_y)$
\sref{IV.5.7.2}. Substituindo $Y$ por $Y'$ e
$X$ por $V\times_Y Y'$, reduzimo-nos, levando
em conta \sref[I]{I.3.6.5}, ao caso em que
$\sh{G}$ satisfaz $(S_k)$ em todo $Y$ e em que,
para todo $y\in f(X)$, $\sh{F}_y$ possui a
propriedade $(S_k)$ em todos os pontos de
$f^{-1}(y)$. Como $\sh{F}$ é $f$-plano,
sabe-se então por \sref{IV.6.4.1} que
$\sh{F}\otimes_{\sh{O}_Y}\sh{G}$ possui a
propriedade $(S_k)$ em $X$, o que demonstra~(i).
Para demonstrar~(ii), basta observar que, sob as
hipóteses feitas, $\sh{G}$ possui a
\oldpage[IV-3]{178}
propriedade $(S_k)$ em uma vizinhança $U$ de $y$
em $Y$ \sref{IV.6.11.2} e
\sref{IV.7.8.3},~(iv). Conclui-se então do mesmo
modo, substituindo desta vez $Y$ por $U$ e $X$
por $V\times_Y U$.
\end{proof}

\begin{theorem}[12.1.6]
\label{IV.12.1.6}
Sejam $f:X\to Y$ um morfismo plano e localmente de
apresentação finita, e $k$ um inteiro. Os seguintes
subconjuntos de $X$ são abertos:
\begin{enumerate}[label=(\roman*)]
\item O conjunto dos pontos $x\in X$ tais que
$X_{f(x)}$ satisfaz a propriedade $(S_k)$ no ponto
$x$.
\item O conjunto dos pontos $x\in X$ tais que
$X_{f(x)}$ satisfaz a propriedade geométrica
$(R_k)$ no ponto $x$, é equidimensional nesse ponto
e, além disso, $x$ não pertence a nenhum ciclo
primo imerso associado a $X_{f(x)}$.
\item O conjunto dos pontos $x\in X$ tais que
$X_{f(x)}$ é geometricamente regular (isto é, liso)
no ponto $x$.
\item O conjunto dos pontos $x\in X$ tais que
$X_{f(x)}$ é geometricamente normal no ponto $x$.
\end{enumerate}
\end{theorem}

\begin{proof}
As etapas A) e B) de \sref{IV.12.0.2} são feitas aqui
como em \sref{IV.12.1.1}: na etapa A), utilizam-se
\sref{IV.6.7.8} e, para~(ii), também
\sref{IV.4.2.7}; na etapa B), utilizam-se
\sref{IV.9.9.2} e \sref{IV.9.9.4}. Resta examinar
a etapa C) em cada caso.
\begin{enumerate}[label=(\roman*)]
\item[(i)] Como em \sref{IV.12.1.1.2}, reduzimo-nos,
usando \sref{IV.6.7.8}, ao caso em que $Y$ é o
espectro de um anel de valorização discreta $A$,
$X=\operatorname{Spec}(B)$, onde $B$ é uma
$A$-álgebra de tipo finito, e $x,x'$ são dois pontos
de $X$ tais que $f(x)=y$ é o ponto fechado e
$y'=f(x')$ é o ponto genérico de $Y$, sendo ainda
$x'$ uma generização de $x$. Como se trata de
demonstrar que $X$ satisfaz $(S_k)$ no ponto $x'$,
podemos substituir $X$ por
$\operatorname{Spec}(\sh{O}_{X,x})$, isto é, supor
que o anel $B$ é local, que o homomorfismo
$A\to B$ é local e que $B$ é uma $A$-álgebra
essencialmente de tipo finito \sref{IV.1.3.8}.
Então, por hipótese, se $t$ é um uniformizante
de $A$, $t$ é um elemento $B$-regular pela
planura, e $B/tB$ satisfaz $(S_k)$. Mas, como
$A$ é um anel universalmente catenário
\sref{IV.5.6.4}, $B$ é catenário; portanto,
por \sref{IV.5.12.4}, $B$ satisfaz $(S_k)$,
o que conclui a demonstração no caso~(i).

\item[(iii)] Aqui a etapa C) é desnecessária:
como $f$ é plano, sabe-se por \sref{IV.6.8.7}
que, quando $Y$ é localmente noetheriano e $f$
é localmente de tipo finito, o conjunto dos
$x\in X$ tais que $X_{f(x)}$ é geometricamente
regular no ponto $x$ é aberto em $X$.

\item[(ii)] Raciocinando como em
\sref{IV.12.1.1.3}, para demonstrar que a
propriedade considerada é estável por
generização, podemos primeiro, considerando
uma extensão finita qualquer $K'$ de $k(y')$
e levando em conta a definição
\sref{IV.6.7.6} da propriedade geométrica
$(R_k)$, bem como a invariância, sob mudança
do corpo de base, das outras duas propriedades
que figuram em~(ii), substituir ali a
propriedade geométrica $(R_k)$ pela propriedade
$(R_k)$. Procedendo então como em~(i),
reduzimo-nos ao caso em que $Y$ é o espectro
de um anel de valorização discreta $A$; como
$A$ é catenário, basta aplicar
\sref{IV.5.12.5} para concluir.

\item[(iv)] O conjunto dos pontos $x\in X$ tais
que $X_{f(x)}$ é geometricamente normal no
ponto $x$ está contido no conjunto dos
$x\in X$ tais que $X_{f(x)}$ é geometricamente
pontualmente integral e satisfaz $(S_2)$ no
ponto $x$; este último é aberto em $X$ por
\sref{IV.12.1.1},~(viii). Podemos, portanto,
supor desde logo que, para todo $x\in X$,
$X_{f(x)}$ é geometricamente pontualmente
integral e satisfaz $(S_2)$ no ponto $x$ e,
\emph{a fortiori}, é equidimensional. Por outro
lado, pelo critério de Serre \sref{IV.5.8.6},
dizer que $X_{f(x)}$ é geometricamente normal
no ponto $x$ significa que $X_{f(x)}$ satisfaz
$(S_2)$ e a propriedade geométrica $(R_1)$
em $x$;
\oldpage[IV-3]{179}
mas, por~(ii) e pelas observações precedentes,
esse conjunto é a interseção de dois abertos
de $X$ e, portanto, é aberto em $X$. C.Q.F.D.
\end{enumerate}
\end{proof}

\begin{corollary}[12.1.7]
\label{IV.12.1.7}
Seja $f:X\to Y$ um morfismo localmente de
apresentação finita. Então o conjunto dos pontos
$x\in X$ nos quais $f$ possui uma das seguintes
propriedades \sref{IV.6.8.1}:
\begin{enumerate}[label=(\roman*)]
\item satisfazer a propriedade $(S_k)$;
\item ter coprofundidade $\leq n$;
\item ser de Cohen--Macaulay;
\item ser regular (ou liso, o que equivale);
\item ser normal;
\item ser reduzido;
\end{enumerate}
é aberto em $X$.
\end{corollary}

\begin{proof}
Com efeito, decorre de \sref{IV.11.3.1} que o
conjunto dos $x\in X$ nos quais $f$ é plano é
aberto. Podemos, portanto, limitar-nos ao caso
em que $f$ é plano; então o corolário decorre
de \sref{IV.12.1.1},~(iv),~(vi) e~(vii), e de
\sref{IV.12.1.6},~(i),~(ii) e~(iv).
\end{proof}

\begin{remarks}[12.1.8]
\label{IV.12.1.8}
\begin{enumerate}[label=(\roman*)]
\item Em \sref{IV.12.1.6},~(ii), não se pode
suprimir a hipótese de que $x$ não pertence
a nenhum ciclo primo imerso associado a
$X_{f(x)}$. Isso é mostrado pelo exemplo
\sref{IV.5.12.3}, em que se toma
$X=\operatorname{Spec}(A)$, sendo $Y$ o espectro
do anel local $A_0$ de $k[T]$ correspondente
ao ideal primo $(T)$, e o morfismo $X\to Y$
correspondendo ao homomorfismo injetivo
$A_0\to A$ obtido, por localização e passagem
ao quociente, da injeção $k[T]\to k[T,U,V]$;
esse morfismo é plano, pois $A$ é um
$A_0$-módulo sem torção \sref[0]{0.6.3.4}.
Aqui a fibra $X_y$ sobre o ponto fechado
$y$ de $Y$, igual a
$\operatorname{Spec}(A/tA)$, é irredutível,
tem dimensão~1 e satisfaz a condição
geométrica $(R_0)$, pois o anel local
em seu ponto genérico é um corpo. Em
contrapartida, sobre o ponto genérico
$y'$ de $Y$, a fibra $X_{y'}$ possui duas
componentes irredutíveis de dimensões
respectivas~0 e~1, e a de dimensão~0
não é reduzida; logo $X_{y'}$ não
satisfaz a condição $(R_0)$.

\item Em \sref{IV.12.1.6},~(ii), tampouco
se pode suprimir a hipótese de que
$X_{f(x)}$ é equidimensional no ponto
$x$. Vê-se isso no exemplo \sref{IV.5.12.6},
com $X=\operatorname{Spec}(A)$ e
$Y=\operatorname{Spec}(A_0)$, onde $A_0$
é definido como em~(i); o morfismo
$X\to Y$ provém novamente da injeção
$k[T]\to k[T,U,V,W]$ por localização
e passagem ao quociente, e é plano
porque $A$ é um $A_0$-módulo sem
torção \sref[0]{0.6.3.4}. A fibra
$X_y$ sobre o ponto fechado $y$ de
$Y$ é reduzida (portanto satisfaz
$(S_1)$) e satisfaz $(R_1)$, mas
possui duas componentes irredutíveis
de dimensões~2 e~1, enquanto a
fibra $X_{y'}$ sobre o ponto genérico
$y'$ de $Y$ não satisfaz $(R_1)$.
\end{enumerate}
\end{remarks}

\subsection{Propriedades locais e globais das fibras
de um morfismo próprio, plano e de apresentação finita}
\label{subsection:IV.12.2}

\begin{theorem}[12.2.1]
\label{IV.12.2.1}
Sejam $f:X\to Y$ um morfismo \emph{próprio}
de apresentação finita, $\sh{F}$ um
$\sh{O}_X$-Módulo $f$-plano e de apresentação
finita, $\Phi$ uma parte finita de
$\mathbf{Z}\cup\{\pm\infty\}$ e $k$ um
inteiro. Os seguintes subconjuntos de $Y$
são abertos:
\begin{enumerate}[label=(\roman*)]
\item O conjunto dos $y\in Y$ tais que o
conjunto das dimensões dos ciclos primos
associados a $\sh{F}_y$ está contido em
$\Phi$.
\item O conjunto dos $y\in Y$ tais que o
conjunto das dimensões das componentes
irredutíveis de
$\operatorname{Supp}(\sh{F}_y)$ contém
$\Phi$.
\item O conjunto dos $y\in Y$ tais que
todos os ciclos primos associados a
$\sh{F}_y$ têm uma mesma dimensão,
igual a $k$.
\item O conjunto dos $y\in Y$ tais que
$\sh{F}_y$ não possui ciclo primo imerso
associado e o conjunto das dimensões
das componentes irredutíveis de
$\operatorname{Supp}(\sh{F}_y)$ é igual a
$\Phi$.
\oldpage[IV-3]{180}
\item O conjunto dos $y\in Y$ tais que
$\sh{F}_y$ satisfaz a propriedade
$(S_k)$ e é equidimensional em todos
os pontos de $X_y$.
\item O conjunto dos $y\in Y$ tais que
$\operatorname{coprof}(\sh{F}_y)\leq k$.
\item O conjunto dos $y\in Y$ tais que
$\sh{F}_y$ é um
$\sh{O}_{X_y}$-Módulo de
Cohen--Macaulay.
\item O conjunto dos $y\in Y$ tais que
$\sh{F}_y$ é geometricamente reduzido.
\item O conjunto dos $y\in Y$ tais que
$\sh{F}_y$ é geometricamente
pontualmente integral em cada ponto
de $X_y$.
\item O conjunto dos $y\in Y$ tais que
$\sh{F}_y$ é geometricamente integral.
\item O conjunto dos $y\in Y$ tais que
$\sh{F}_y$ não possui ciclo primo
imerso associado e a soma $l(y)$ das
multiplicidades totais
\sref{IV.4.7.12} de $\sh{F}_y$ nos
pontos maximais de
$\operatorname{Supp}(\sh{F}_y)$ é
$\leq k$.
\end{enumerate}
\end{theorem}

\begin{proof}
Com exceção de~(ii), (iii), (iv), (x)
e~(xi), as propriedades $P(y)$
consideradas têm a forma: “para todo
$x\in X_y$, $\sh{F}_y$ possui a
propriedade $Q(x)$”, em que $Q(x)$
é uma das propriedades (i) a (viii)
de \sref{IV.12.1.1}. Por
\sref{IV.12.1.1}, o conjunto $U$ dos
$x\in X$ para os quais $Q(x)$ é
verdadeira é aberto, e o conjunto
$V$ dos $y\in Y$ para os quais
$P(y)$ é verdadeira é precisamente
$Y-f(X-U)$; nesses casos, portanto,
o teorema já vale supondo apenas que
o morfismo $f$ seja \emph{fechado}.
Para~(iii), aplica-se~(i) com
$\Phi$ reduzido a um único elemento.
Resta examinar separadamente os
casos~(ii), (iv) e~(xi), pois~(x)
decorre imediatamente de~(xi) e de
\sref{IV.4.7.13}\footnote{A fonte
impressa cita (4.7.14), mas a Seção
IV.4.7 termina com a Proposição
4.7.13, antes do início de IV.4.8;
4.7.13 é o critério pertinente de
integralidade geométrica.},
aplicando sempre o método descrito
em \sref{IV.12.0.2}.
\end{proof}

\begin{env}[12.2.1.1]
\label{IV.12.2.1.1}
Caso~(ii): As etapas A) e B) da
demonstração procedem como no início de
\sref{IV.12.1.1}; para a etapa A),
utilizam-se \sref{IV.8.9.1},
\sref{IV.11.2.6} e \sref{IV.4.2.7};
para a etapa B), utiliza-se
\sref{IV.9.5.5} aplicado a
$\operatorname{Supp}(\sh{F})$. Resta
demonstrar que, se (supondo $Y$
noetheriano) o conjunto das dimensões
das componentes irredutíveis de
$\operatorname{Supp}(\sh{F}_y)$ contém
$\Phi$, então o mesmo vale para o
conjunto das dimensões das componentes
irredutíveis de
$\operatorname{Supp}(\sh{F}_{y'})$,
para toda \emph{generização} $y'$
de $y$ em $Y$. Seja $Y_1$ o espectro
de um anel de valorização discreta
tal que, sendo $y_1$ e $y'_1$ o ponto
fechado e o ponto genérico de $Y_1$,
exista um morfismo $h:Y_1\to Y$ com
$h(y_1)=y$ e $h(y'_1)=y'$
\sref[II]{II.7.1.7}. Aplicando
\sref{IV.4.2.7}, vemos que podemos
substituir $Y$ por $Y_1$ e $X$ por
$X_1=X\times_Y Y_1$; em outras
palavras, basta considerar o caso em
que $Y$ é o espectro de um anel de
valorização discreta, $y$ é seu
ponto fechado e $y'$ seu ponto
genérico.

Usando (\textbf{Err}\textsubscript{\textbf{III}},~30),
podemos restringir-nos ao caso em que
$\operatorname{Supp}(\sh{F})=X$, de
modo que $f$ é quase-plano
\sref{IV.2.3.3}; as componentes
irredutíveis $Z_i$ de $X$ dominam
então $Y$ \sref{IV.2.3.4}, isto é,
seus pontos genéricos pertencem a
$X_{y'}$, e os $Z_i\cap X_{y'}$
são as componentes irredutíveis de
$X_{y'}$. Mas toda componente
irredutível de $X_y$ está contida
em algum dos $Z_i$ e, portanto,
é uma componente irredutível de
$(Z_i)_y$; ora, por
\sref{IV.7.1.13}, as dimensões de
todas as componentes irredutíveis
de $(Z_i)_y$ são iguais a
$\dim((Z_i)_{y'})$, o que conclui
a demonstração de~(ii).
\end{env}

\begin{env}[12.2.1.2]
\label{IV.12.2.1.2}
Caso~(iv): O mesmo raciocínio do
início de \sref{IV.12.2.1.1},
utilizando \sref{IV.12.1.1},~(iii),
mostra que já podemos supor,
substituindo $Y$ por um aberto de
$Y$, que, para todo $y\in Y$,
$\sh{F}_y$ não possui ciclo primo
imerso associado. A afirmação do
caso~(iv) decorre então
imediatamente das afirmações dos
casos~(i) e~(ii).
\end{env}

\begin{env}[12.2.1.3]
\label{IV.12.2.1.3}
\oldpage[IV-3]{181}
Caso~(xi): Para a etapa A) do
raciocínio, utilizam-se
\sref{IV.8.9.1},
\sref{IV.11.2.6} e
\sref{IV.8.10.5},~(xii)
(para conservar a hipótese de
que $f$ é próprio), bem como
\sref{IV.4.2.7},
\sref{IV.4.5.6} e
\sref{IV.4.7.9}. Para a etapa
B), vemos, como no caso~(iv),
que podemos supor que, para
todo $y\in Y$, $\sh{F}_y$ não
possui ciclo primo imerso
associado, e aplicamos
\sref{IV.9.8.8}, que demonstra
ser construtível a função
$z\mapsto l(z)$. Resta,
portanto, mostrar, supondo
$Y$ noetheriano e $f$ próprio,
que, para toda generização
$y'$ de um ponto $y\in Y$,
$l(y')\leq l(y)$. Raciocinando
como em \sref{IV.12.1.1.3},
vemos, com auxílio de
\sref{IV.4.7.8} e
\sref{IV.4.5.11}, que podemos
supor que as componentes
irredutíveis de
$\operatorname{Supp}(\sh{F}_{y'})$
são geometricamente
irredutíveis e que $l(y')$
é a \emph{soma dos
comprimentos} dos
$(\sh{F}_{y'})_{z'_j}$ nos
pontos maximais $z'_j$ de
$\operatorname{Supp}(\sh{F}_{y'})$.
Aplicando novamente
\sref{IV.4.2.7},
\sref{IV.4.5.6} e
\sref{IV.4.7.9},
reduzimo-nos, como em
\sref{IV.12.2.1.1}, ao caso
em que $Y$ é o espectro de um
anel de valorização discreta,
$y$ seu ponto fechado e $y'$
seu ponto genérico. A
desigualdade $l(y')\leq l(y)$
será então consequência do

\begin{lemma}[12.2.1.4]
\label{IV.12.2.1.4}
Sejam $Y$ o espectro de um anel
de valorização discreta, $y$
seu ponto fechado, $y'$ seu
ponto genérico, $f:X\to Y$
um morfismo próprio, $\sh{F}$
um $\sh{O}_X$-Módulo coerente e
$f$-plano, e $z_i$ (resp.\
$z'_j$) os pontos maximais de
$\operatorname{Supp}(\sh{F}_y)$
(resp.\
$\operatorname{Supp}(\sh{F}_{y'})$).
Suponhamos que $\sh{F}_y$ não
possua ciclo primo imerso
associado. Então,
\[
\sum_j\operatorname{long}
((\sh{F}_{y'})_{z'_j})
\leq
\sum_i\operatorname{long}
((\sh{F}_y)_{z_i}).
\]
\end{lemma}

Temos $Y=\operatorname{Spec}(A)$,
onde $A$ é um anel de
valorização discreta, cujo
uniformizante designaremos
por $t$, de modo que
$\sh{F}_y=\sh{F}/t\sh{F}$.
Como $\sh{F}$ é $f$-plano,
os $z'_j$ também são os
pontos maximais de
$\operatorname{Supp}(\sh{F})$
\sref{IV.2.3.4}; para cada
$z_i$, designemos por
$z'_{ik}$ aqueles $z'_j$
que são generizações de
$z_i$; resulta de
\sref{IV.3.4.1.1} que
\[
\operatorname{long}
((\sh{F}_y)_{z_i})
\geq
\sum_k\operatorname{long}
((\sh{F}_{y'})_{z'_{ik}}),
\]
donde, somando,
\[
\sum_i\operatorname{long}
((\sh{F}_y)_{z_i})
\geq
\sum_i\sum_k
\operatorname{long}
((\sh{F}_{y'})_{z'_{ik}}).
\]

O lema estará demonstrado se
estabelecermos que, para todo
$z'_j$, existe ao menos um
índice $i$ tal que $z'_j$
é um dos $z'_{ik}$. Ora,
como $f$ é próprio (portanto,
fechado) e $f(z'_j)=y'$,
existe $x\in X$ que é uma
especialização de $z'_j$ e
satisfaz $f(x)=y$; em outras
palavras,
$\overline{\{z'_j\}}\cap X_y$
não é vazio. Como $t$ é
$\sh{F}$-regular pela
planura, deduzimos de
\sref{IV.3.4.3} que existe
ao menos um ponto de
$\operatorname{Ass}(\sh{F}_y)$
do qual $z'_j$ é uma
generização; mas, por
hipótese, os pontos de
$\operatorname{Ass}(\sh{F}_y)$
são os $z_i$, o que conclui
a demonstração.
\end{env}

\begin{corollary}[12.2.2]
\label{IV.12.2.2}
Sob as hipóteses de
\sref{IV.12.2.1}:
\begin{enumerate}[label=(\roman*)]
\item A função
$y\mapsto\dim(\operatorname{Supp}(\sh{F}_y))$
é contínua (portanto, localmente
constante) em $Y$.
\item A função
$y\mapsto\operatorname{coprof}(\sh{F}_y)$
\sref{IV.5.7.1} é semicontínua
superiormente em $Y$.
\item A função $y\mapsto l(y)$
\sref{IV.12.2.1},~(xi) é
semicontínua superiormente em
$Y$ quando os $\sh{F}_y$ não
possuem ciclo primo imerso
associado.
\oldpage[IV-3]{182}
\item A função
$y\mapsto\operatorname{prof}^{*}(\sh{F}_y)$
\sref{IV.10.8.1} é semicontínua
inferiormente em $Y$.
\end{enumerate}
\end{corollary}

\begin{proof}
\emph{(i)} De \sref{IV.12.2.1},~(i),
aplicado a $\Phi$ igual ao menor
intervalo de $\mathbf{N}$ que
contém as dimensões dos ciclos
primos associados a $\sh{F}_y$,
resulta que
$z\mapsto\dim(\operatorname{Supp}(\sh{F}_z))$
é semicontínua superiormente no
ponto $y$; por outro lado, de
\sref{IV.12.2.1},~(ii), aplicado
a $\Phi$ igual ao conjunto das
dimensões das componentes
irredutíveis de
$\operatorname{Supp}(\sh{F}_y)$,
resulta que a mesma função é
semicontínua inferiormente no
ponto $y$, donde a conclusão.
As afirmações~(ii) e~(iii)
são apenas outras formulações
de \sref{IV.12.2.1},~(vi)
e~(xi). Por fim, segundo
\sref{IV.12.1.3}, o conjunto
dos $x\in X$ tais que
$\operatorname{prof}_{f(x)}^{*}
(\sh{F}_{f(x)})\leq k$ é aberto,
e a conclusão de~(iv) decorre
do mesmo raciocínio empregado
no início de \sref{IV.12.2.1}.
\end{proof}

\begin{remarks}[12.2.3]
\label{IV.12.2.3}
\begin{enumerate}[label=(\roman*)]
\item Observe-se que, em
\sref{IV.12.2.1},~(ii), podemos
dispensar a hipótese de que $f$
seja \emph{próprio}.

\item Em \sref{IV.12.2.1},~(ii),
não podemos substituir a condição
de que o conjunto $E(y)$ das
dimensões das componentes
irredutíveis de
$\operatorname{Supp}(\sh{F}_y)$
contenha $\Phi$ pela condição de
que $E(y)$ seja \emph{igual a}
$\Phi$. Com efeito, seja $k$ um
corpo e consideremos o espaço
projetivo tridimensional
$P=\mathbf{P}^3_k=
\operatorname{Proj}(C)$, onde
$C=k[t_0,t_1,t_2,t_3]$ (com os
$t_i$ indeterminadas). Em $P$,
seja $X_0$ o subpré-esquema
fechado “reunião da reta $X_1$
definida por $t_1=t_3=0$ e do
plano $X_2$ definido por
$t_2-t_3=0$”. Isso expressa
geometricamente que
$X_0=\operatorname{Proj}
(C/\mathfrak a)$, onde o ideal
graduado $\mathfrak a$ é
$\mathfrak b\mathfrak c$, com
$\mathfrak b=Ct_1+Ct_3$ e
$\mathfrak c=C(t_2-t_3)$.
Consideremos o morfismo
$f_0:X_0\to X_1$ que,
geometricamente, é a “projeção
de $X_0$ sobre $X_1$ a partir
da reta no infinito $D$ definida
por $t_0=t_2=0$”. Em termos
algébricos, $f_0$ corresponde
ao homomorfismo de álgebras
graduadas
\[
k[t_0,t_2]\longrightarrow
k[t_0,t_1,t_2,t_3]/\mathfrak a
\]
obtido por passagem ao quociente
a partir da injeção canônica
$k[t_0,t_2]\to
k[t_0,t_1,t_2,t_3]$. É claro
que $f_0$ é um morfismo
projetivo, portanto
\emph{próprio}; o mesmo vale
para sua restrição $f:X\to Y$,
onde $Y=D_+(t_0)$ e
$X=f_0^{-1}(Y)$. Para ver que
$f$ é \emph{plano}, basta
observar que $k[t_2]$ é um
domínio de ideais principais
e que o anel
\[
C'=k[t_1,t_2,t_3]/
(t_2-t_3)((t_1)+(t_3))
\]
é um $k[t_2]$-módulo sem
torção \sref[0]{0.6.3.4}.
Para todo $y\in Y$ diferente
do ponto $y_0$ definido por
$t_2=0$, $f^{-1}(y)$ tem,
portanto, duas componentes
irredutíveis, de dimensões~0
e~1, enquanto $f^{-1}(y_0)$
possui uma só componente
irredutível, de dimensão~1.

\item Em \sref{IV.12.2.1},~(i),
tampouco podemos substituir a
condição de que o conjunto
$E'(y)\supset E(y)$ das
dimensões dos ciclos primos
associados a $\sh{F}_y$
esteja contido em $\Phi$
pela condição de que seja
\emph{igual a} $\Phi$. Sejam
$k$ um corpo, $A$ o anel
de polinômios $k[t]$ e $B$
o subanel
$k[t^4,t^3u,tu^3,u^4]$
de $k[t,u]$ ($t,u$
indeterminadas). O anel $B$
não é integralmente fechado:
$t^2u^2$ pertence a seu
fecho integral, mas não a
$B$. Se $\mathfrak m$ é o
ideal maximal de $B$
gerado por $t^4,t^3u,tu^3$
e $u^4$, então $\mathfrak m$
é um ideal primo associado
imerso de $B/Bt^4$%
\footnote{A fonte impressa
traz $A/At^4$; a construção
e o cálculo da fibra que
seguem exigem $B/Bt^4$.}.
Ponhamos
$Y=\operatorname{Spec}(A)$,
$X=\operatorname{Spec}(B)$,
e seja $f:X\to Y$ o
morfismo correspondente ao
homomorfismo $A\to B$ que
envia $t$ em $t^4$. Como
esse homomorfismo faz de
$B$ um $A$-módulo sem
torção, $f$ é plano
\sref[0]{0.6.3.4}. Se $y$
é o ponto de $Y$
correspondente ao ideal
maximal $tA$, o pré-esquema
$f^{-1}(y)$ é irredutível
e de dimensão~1, mas possui
um ciclo primo imerso
associado de dimensão~0.
Em contrapartida, se $y'$
é o ponto genérico de $Y$,
$f^{-1}(y')$ não tem ciclo
primo imerso associado,
pois é o espectro de uma
$k(y')$-álgebra integral
de dimensão~1. Neste
exemplo, $f$ é um
morfismo afim não próprio;
obtemos imediatamente um
exemplo análogo em que $f$
é próprio e plano,
considerando $X$ como
um aberto denso em toda
parte de um esquema
projetivo sobre $Y$
\sref[II]{II.5.3.4} e
\sref[II]{II.5.3.2},
ou procedendo diretamente
como na Observação~(ii).

As Observações~(ii) e~(iii)
mostram a necessidade de
incluir em
\sref{IV.12.2.1},~(iv)
a condição de que
$\sh{F}_y$ não possua
ciclo primo imerso
associado.

\item Em \sref{IV.12.2.1},~(xi),
não podemos suprimir a hipótese
de que $f$ seja próprio.
Com efeito, sejam $k$ um corpo,
$Y$ a “reta afim”, espectro
do anel de polinômios $k[t]$
($t$ indeterminada), e $X$
a soma disjunta de $Y$ com
o aberto complementar do
ponto fechado $y$ de $Y$
definido por $t=0$. É claro
que o morfismo $f:X\to Y$,
igual às injeções canônicas
nas duas componentes de $X$,
é plano e que, tomando
$\sh{F}=\sh{O}_X$, nenhum
$\sh{F}_z$ possui ciclo
primo imerso associado
para $z\in Y$. Contudo,
$l(y)=1$ e $l(z)=2$
para todo $z\ne y$.

\item Em \sref{IV.12.2.1},~(xi),
tampouco podemos suprimir a
hipótese de que $\sh{F}_y$
não possua ciclo primo
imerso associado. Isso
é mostrado pelo exemplo
da Observação~(ii): com
$\sh{F}=\sh{O}_X$,
$l(y_0)=1$ e $l(y)=2$
para $y\ne y_0$ em $Y$.

\item Não sabemos se, em
\sref{IV.12.2.1},~(xi),
podemos substituir a
hipótese de que $\sh{F}$
é $f$-plano pela de que
$\operatorname{Supp}(\sh{F})=X$
e $f$ é universalmente
aberto. Retomando a
demonstração de
\sref{IV.12.2.1.4},
vemos, usando também
\sref{IV.14.3.6}, que
seria necessário resolver
a seguinte questão: sejam
$A$ um anel local
noetheriano, $M$ um
$A$-módulo de tipo finito,
\oldpage[IV-3]{183}
$t$ um elemento do ideal
maximal $\mathfrak m$
de $A$ que não está
contido em nenhum ideal
primo minimal de $A$;
se existe um desses
ideais primos minimais
$\mathfrak p$ tal que
$\dim(A/\mathfrak p)=1$,
é verdade que
$\mathfrak m$ é um ideal
primo associado a
$M/tM$ (sem supor que
$t$ seja $M$-regular)?

Observe-se, contudo, que
em \sref{IV.12.2.1},~(xi)
não podemos simplesmente
suprimir a hipótese de que
$\sh{F}$ é $f$-plano.
Tomemos aqui o $P$ da
Observação~(ii) e, em $P$,
o subpré-esquema fechado
$X_0$ que é a reunião
das três retas
$X_1,X_2,X_3$ definidas,
respectivamente, por
$t_1=t_3=0$,
$t_1-t_0=t_3=0$ e
$t_2=t_3=0$. Definamos
a projeção $f_0$ de
$X_0$ sobre $X_1$ como
antes, e sua restrição
$f:X\to Y$, onde
$Y=D_+(t_0)$ é a reta
afim e $X=f_0^{-1}(Y)$.
O morfismo $f$ é próprio,
mas não plano, e, tomando
$\sh{F}=\sh{O}_X$,
$\sh{F}_y$ não possui
ciclo primo imerso
associado para nenhum
$y\in Y$. Mas
$l(y_0)=1$ e $l(y)=2$
para $y\ne y_0$ em $Y$.
\end{enumerate}
\end{remarks}

\begin{theorem}[12.2.4]
\label{IV.12.2.4}
Sejam $f:X\to Y$ um morfismo
próprio, plano e de apresentação
finita, e $k$ um inteiro
$\geqslant 1$. Os seguintes
subconjuntos de $Y$ são abertos:
\begin{enumerate}[label=(\roman*)]
\item O conjunto dos $y\in Y$
tais que $X_y$ possui a
propriedade $(\mathrm{S}_k)$.
\item O conjunto dos $y\in Y$
tais que $X_y$ satisfaz a
propriedade geométrica
$(\mathrm{R}_k)$, é
equidimensional em cada ponto
e não possui ciclo primo
imerso associado.
\item O conjunto dos $y\in Y$
tais que $X_y$ é
geometricamente regular
(isto é, suave sobre $k(y)$).
\item O conjunto dos $y\in Y$
tais que $X_y$ é
geometricamente normal.
\item O conjunto dos $y\in Y$
tais que $X_y$ é
geometricamente reduzido.
\item O conjunto dos $y\in Y$
tais que $X_y$ é
geometricamente reduzido e
o número geométrico de
componentes conexas de $X_y$
é igual a $k$.
\item O conjunto dos $y\in Y$
tais que $X_y$ é
geometricamente pontualmente
integral \sref{IV.4.6.9}.
\item O conjunto dos $y\in Y$
tais que $X_y$ é
geometricamente integral.
\item O conjunto dos $y\in Y$
tais que $X_y$ não possui
ciclo primo imerso associado
e a multiplicidade total
\sref{IV.4.7.4} de $X_y$
sobre $k(y)$ é $\leqslant k$.
\end{enumerate}
\end{theorem}

\begin{proof}
Com exceção de~(vi), essas
afirmações são casos particulares
de afirmações de
\sref{IV.12.2.1}, ou se deduzem
de afirmações de
\sref{IV.12.1.6} como no início
da demonstração de
\sref{IV.12.2.1}. Para~(vi),
reduzimo-nos, como sempre
(levando em conta a invariância
do número geométrico de
componentes conexas por mudança
de base \sref{IV.4.5.6}),
ao caso em que $Y$ é noetheriano.
O conjunto $U$ dos $y\in Y$
tais que $X_y$ é geometricamente
reduzido é aberto em $Y$, em
virtude de~(v). De
\sref[III]{III.7.8.7} e
\sref[III]{III.7.8.6} resulta
então que, para todo $y\in U$,
existe uma vizinhança
$V\subset U$ de $y$ e um
inteiro $m$ tais que, para
\emph{todo} $z\in V$,
$\Gamma(X_z,\sh{O}_{X_z})$
é isomorfo a $(k(z))^m$.
Mas, em virtude de
\sref[III]{III.4.3.4}, $m$
é então o número geométrico
de componentes conexas de
$X_z$, de onde a conclusão.
Daremos outra demonstração
de~(vi) em \sref{IV.15.5.9}.
\end{proof}

\subsection{Propriedades cohomológicas locais das fibras de um morfismo plano e localmente de apresentação finita}
\label{subsection:IV.12.3}

\begin{lemma}[12.3.1]
\label{IV.12.3.1}
Sejam $A$ um anel, $B$ uma
$A$-álgebra de apresentação
finita que é um $A$-módulo
plano, e $M$ um $B$-módulo
de apresentação finita que
é um $A$-módulo plano.
Então, as seguintes condições
são equivalentes:
\begin{enumerate}[label=\emph{\alph*)}]
\item $M$ é um $B$-módulo
projetivo.
\item $M$ é um $B$-módulo
plano.
\item Para todo $y\in\Spec(A)$,
$M\otimes_A k(y)$ é um
$(B\otimes_A k(y))$-módulo
projetivo.
\end{enumerate}
\end{lemma}

\begin{proof}
\oldpage[IV-3]{184}
A equivalência de \emph{a)}
e \emph{b)} resulta de
Bourbaki, \emph{Alg.\ comm.},
cap.~II, \textsection\,5,
n\textsuperscript{o}\,2,
cor.~2 do teo.~1.
Como \emph{a)} implica
trivialmente \emph{c)},
resta provar que \emph{c)}
implica \emph{b)}.
Isso resulta do critério
de planicidade por fibras
\sref{IV.11.3.10}, aplicado
com $g=h$, $f=1_X$.
\end{proof}

\begin{proposition}[12.3.2]
\label{IV.12.3.2}
Sejam $A$ um anel, $B$ uma
$A$-álgebra de apresentação
finita que é um $A$-módulo
plano, e $M$ um $B$-módulo
de apresentação finita que
é um $A$-módulo plano.
Então:
\begin{enumerate}[label=(\roman*)]
\item Existe uma resolução
à esquerda de $M$ por
$B$-módulos livres de
tipo finito.
\item Temos
\begin{equation}
\tag{12.3.2.1}
\label{IV.12.3.2.1}
\operatorname{dim.proj}_B(M)
=\sup_{y\in\Spec(A)}
\operatorname{dim.proj}_{B\otimes_A k(y)}(M\otimes_A k(y))
=\operatorname{Tor.dim}_B(M).
\end{equation}
Aqui, $\operatorname{Tor.dim}_B(M)$
é o \emph{menor inteiro} $i$
tal que
$\operatorname{Tor}_j^B(M,N)=0$
para todo $j>i$ e todo
$B$-módulo $N$
(e é $+\infty$ se não
existe tal inteiro $i$).\footnote{A impressão
e os testemunhos francês, espanhol
e inglês alinhados trazem $j\geqslant i$.
A leitura corrente $j>i$ é uma
correção editorial provisória, encaminhada
ao responsável pelo cânone. A desigualdade
estrita é necessária à igualdade com a
dimensão projetiva em \textup{(12.3.2.1)}:
um módulo livre não nulo tem dimensão
projetiva zero, mas seu $\operatorname{Tor}_0$
não é identicamente nulo. A leitura
impressa fica preservada nesta nota.}
\end{enumerate}
\end{proposition}

\begin{proof}
\begin{enumerate}[label=(\roman*)]
\item Em virtude de
\sref{IV.8.9.1}, existem
um subanel noetheriano $A_0$
de $A$, uma $A_0$-álgebra
de tipo finito $B_0$ e um
$B_0$-módulo de tipo finito
$M_0$ tais que
$B\simeq B_0\otimes_{A_0}A$
e $M\simeq M_0\otimes_{A_0}A$.
Por \sref{IV.11.2.7}, podemos
supor, além disso, que $B_0$
e $M_0$ são $A_0$-módulos
planos. Existe então uma
resolução à esquerda
$(L_0)_\bullet$ de $M_0$
formada por $B_0$-módulos
livres de tipo finito.
Como $M_0$ e os $(L_0)_i$
são $A_0$-módulos planos,
$(L_0)_\bullet\otimes_{A_0}A$
é uma resolução à esquerda
de $M$ formada por
$B$-módulos livres de
tipo finito
\sref{IV.2.1.10}.

\item Em virtude de
\sref[0]{0.17.2.2},~(ii),
$\operatorname{Tor.dim}_B(M)\leqslant\operatorname{dim.proj}_B(M)$,
e a definição de dimensão
projetiva dá imediatamente
\[
\operatorname{dim.proj}_{B\otimes_A k(y)}(M\otimes_A k(y))
\leqslant\operatorname{dim.proj}_B(M)
\]
para todo $y\in\Spec(A)$.
Para provar as desigualdades
inversas, consideremos uma
resolução à esquerda $(L_i)$
de $M$ por $B$-módulos livres
de tipo finito, e suponhamos
que $\operatorname{Tor.dim}_B(M)=n$
(respectivamente,
$\operatorname{dim.proj}_{B\otimes_A k(y)}(M\otimes_A k(y))\leqslant n$
para todo $y\in\Spec(A)$).
Então,
$R=\operatorname{Im}(L_n\to L_{n-1})=\operatorname{Ker}(L_{n-1}\to L_{n-2})$
é um $B$-módulo de tipo
finito e também um $A$-módulo
plano, em virtude das hipóteses
sobre $M$ e $B$ e de
\sref{IV.2.1.10}.
Além disso, temos
\[
\operatorname{Tor}_{n+1}^B(M,N)=\operatorname{Tor}_1^B(R,N)
\]
para todo $B$-módulo $N$
(M, V, 7). A hipótese
$\operatorname{Tor.dim}_B(M)=n$
implica, portanto,
$\operatorname{Tor}_1^B(R,N)=0$
para todo $B$-módulo $N$;
isto é, $R$ é um $B$-módulo
plano, logo projetivo em
virtude de \sref{IV.12.3.1}.
Isso estabelece que
$\operatorname{dim.proj}_B(M)\leqslant n$.
Por outro lado, a hipótese
$\operatorname{dim.proj}_{B\otimes_A k(y)}(M\otimes_A k(y))\leqslant n$
para todo $y\in\Spec(A)$
implica, por tensorização
com $k(y)$, que em cada
uma das sequências seguintes
(exatas em virtude da
planicidade sobre $A$ de
$M$, dos $L_i$ e de $R$
\sref{IV.2.1.10})
\[
0\to R\otimes_A k(y)\to L_{n-1}\otimes_A k(y)
\to\cdots\to L_0\otimes_A k(y)\to M\otimes_A k(y)\to0
\]
o $(B\otimes_A k(y))$-módulo
$R\otimes_A k(y)$ é
projetivo, para todo
$y\in\Spec(A)$.
Concluímos então, novamente
por \sref{IV.12.3.1}, que
$R$ é um $B$-módulo
projetivo, de modo que
$\operatorname{dim.proj}_B(M)\leqslant n$.
\end{enumerate}
\end{proof}

\begin{proposition}[12.3.3]
\label{IV.12.3.3}
Sejam $f:X\to S$ um morfismo localmente de apresentação finita e
$\sh{L}_\bullet$ um complexo de $\sh{O}_X$-Módulos quase-coerentes
de apresentação finita. Para todo $s\in S$, denotemos por
$(\sh{L}_\bullet)_s$ o complexo
$\sh{L}_\bullet\otimes_{\sh{O}_S}k(s)$ de
$\sh{O}_{X_s}$-Módulos de tipo finito.
Suponhamos que $\sh{L}_{n-1}$ seja $f$-plano.
Então, o conjunto $U$ dos $x\in X$ tais que
$(\sh{H}_n((\sh{L}_\bullet)_{f(x)}))_x=0$ é aberto em $X$.
Se, além disso, $\sh{L}_n$ é $f$-plano, temos
$\sh{H}_n(\sh{L}_\bullet)|_U=0$.
\end{proposition}

\begin{proof}
É claro que basta considerar o caso $n=1$, no qual o complexo
$\sh{L}_\bullet$ se reduz a
$0\to\sh{L}_2\xrightarrow{u}\sh{L}_1\xrightarrow{v}\sh{L}_0\to0$.
Podemos primeiro reduzir ao caso em que $S=\Spec(A)$ e $X$
\oldpage[IV-3]{185}
são afins, e depois ao caso em que $S$ é \emph{noetheriano}.
Com efeito, por \sref{IV.8.9.1} e \sref{IV.8.5.2}, existem
um pré-esquema noetheriano $S'=\Spec(A')$, onde $A'$ é um subanel
de $A$, um morfismo de tipo finito $f':X'\to S'$ e três
$\sh{O}_{X'}$-Módulos coerentes $\sh{L}'_i$
($0\leqslant i\leqslant2$), tais que
$X=X'\times_{S'}S$, $f=f'_{(S)}$ e
$\sh{L}_i=\sh{L}'_i\otimes_{S'}S$, juntamente com dois
homomorfismos $u'$, $v'$ tais que $u=u'\otimes1$,
$v=v'\otimes1$ e $v'\circ u'=0$.
Podemos supor, além disso, que $\sh{L}'_0$ é $f'$-plano
\sref{IV.11.2.7} e que $\sh{L}'_1$ é $f'$-plano quando
se supõe que $\sh{L}_1$ seja $f$-plano.
Por planicidade fiel, a hipótese
$(\sh{H}_n((\sh{L}_\bullet)_{f(x)}))_x=0$
equivale a
$(\sh{H}_n((\sh{L}'_\bullet)_{f'(x')}))_{x'}=0$,
onde $x'$ é a projeção de $x$ em $X'$.
Assim, se $U'$ é o conjunto dos $x'\in X'$ tais que
$(\sh{H}_n((\sh{L}'_\bullet)_{f'(x')}))_{x'}=0$,
$U$ é a imagem inversa de $U'$ pela projeção $X\to X'$.
Isso reduz a primeira afirmação ao caso em que $S$ é noetheriano.

Para a segunda afirmação, observe-se também que $A$ é o limite
indutivo filtrante de seus subanéis $A_\lambda$ que são
$A'$-álgebras de tipo finito. Ponhamos $S_\lambda=\Spec(A_\lambda)$,
$X_\lambda=X'\times_{S'}S_\lambda$ e
$\sh{L}^{(\lambda)}_i=\sh{L}'_i\otimes_{S'}S_\lambda$; sejam
$u_\lambda=u'\otimes1:\sh{L}_2^{(\lambda)}\to\sh{L}_1^{(\lambda)}$ e
$v_\lambda=v'\otimes1:\sh{L}_1^{(\lambda)}\to\sh{L}_0^{(\lambda)}$,
de modo que
$\sh{H}_1(\sh{L}_\bullet^{(\lambda)})=\operatorname{Ker}(v_\lambda)/\operatorname{Im}(u_\lambda)$.
Como os limites indutivos filtrantes são exatos na categoria
dos grupos abelianos, temos
\[
\operatorname{Ker}(v)=\varinjlim_\lambda\operatorname{Ker}(v_\lambda),\qquad
\operatorname{Im}(u)=\varinjlim_\lambda\operatorname{Im}(u_\lambda),\qquad
\sh{H}_1(\sh{L}_\bullet)=\varinjlim_\lambda\sh{H}_1(\sh{L}_\bullet^{(\lambda)}).
\]
Se tivermos suposto que $\sh{L}_1$ é $f$-plano e, reduzindo
ao caso $X=U$, tivermos provado que
$\sh{H}_1(\sh{L}_\bullet^{(\lambda)})=0$ para todo $\lambda$,
deduziremos, de fato, que $\sh{H}_1(\sh{L}_\bullet)=0$.

\medskip
I) Suponhamos, de agora em diante, que $S$ seja noetheriano.
Sabemos, sem nenhuma hipótese de planicidade sobre $\sh{L}_0$,
que o conjunto $U$ é \emph{construtível} em $X$
\sref{IV.9.9.6}. Usando agora
\hyperref[0.9.2.5]{($\mathbf{0}_{\mathrm{III}}$,~9.2.5)},
resta mostrar que toda \emph{generização} $x'$ em $X$ de
um ponto $x\in U$ também pertence a $U$.
O método exposto em \sref{IV.12.0.2} aplica-se sem alteração,
levando em conta que, para um pré-esquema $Z$ sobre um corpo
$k$ e uma extensão $k'$ de $k$, a projeção
$Z\otimes_k k'\to Z$ é fielmente plana.
Podemos, portanto, supor que $S$ seja o espectro de um anel
de valuação discreta $A$, com uniformizante $t$, que $x$ esteja
sobre o ponto fechado de $S$ e que $x'$ esteja sobre o ponto
genérico de $S$.
A hipótese de que $\sh{L}_0$ é $f$-plano implica que $t$ é
$\sh{L}_0$-regular
\hyperref[0.6.3.4]{($\mathbf{0}_{\mathrm{I}}$,~6.3.4)}.
Reduzimo-nos, assim, à prova do lema seguinte, no qual
$B$, $M$, $N$ e $P$ serão substituídos, respectivamente, por
$\sh{O}_x$, $(\sh{L}_2)_x$, $(\sh{L}_1)_x$ e $(\sh{L}_0)_x$.

\begin{lemma}[12.3.3.1]
\label{IV.12.3.3.1}
Sejam $B$ um anel local noetheriano, $M$, $N$ e $P$ três
$B$-módulos, com $N$ de tipo finito, e
$M\xrightarrow{u}N\xrightarrow{v}P$ dois homomorfismos
tais que $v\circ u=0$.
Seja $t$ um elemento do ideal maximal de $B$ que é $P$-regular.
Se a sequência
\begin{equation}
\tag{12.3.3.2}
\label{IV.12.3.3.2}
M/tM\xrightarrow{u\otimes1_{B/tB}}N/tN
\xrightarrow{v\otimes1_{B/tB}}P/tP
\end{equation}
é exata, então a sequência $M\xrightarrow{u}N\xrightarrow{v}P$
é exata.
\end{lemma}

Observemos primeiro que, se substituirmos $M$ por sua imagem
$Q$ em $N$ e $u$ pela injeção $j:Q\to N$, a imagem de
$j\otimes1$ é a mesma que a de $u\otimes1$.
Podemos, portanto, sem alterar nem a hipótese nem a conclusão,
supor que $u$ seja injetivo.
Por outro lado, se $R$ é a imagem de $v$ e $\rho:N\to R$
é a sobrejeção canônica, $t$ é evidentemente $R$-regular,
$\rho\circ u=0$, e o núcleo de $\rho\otimes1$ está contido
no de $v\otimes1$.
Logo, é igual a ele se a sequência \sref{IV.12.3.3.2} é
exata. Como, além disso,
$\operatorname{Ker}(\rho)=\operatorname{Ker}(v)$,
\oldpage[IV-3]{186}
vemos que, para provar o lema, também podemos supor que
$v$ seja sobrejetivo. O lema será então consequência do seguinte.

\begin{lemma}[12.3.3.3]
\label{IV.12.3.3.3}
Sejam $B$ um anel, $M$, $N$ e $P$ três $B$-módulos, e
$u:M\to N$, $v:N\to P$ dois homomorfismos tais que $u$
seja injetivo, $v$ seja sobrejetivo e $v\circ u=0$.
Seja $t$ um elemento de $B$ que é $P$-regular.
Então, a menos de um isomorfismo canônico, temos
\begin{equation}
\tag{12.3.3.4}
\label{IV.12.3.3.4}
\operatorname{Ker}(v\otimes1)/\operatorname{Im}(u\otimes1)
=(\operatorname{Ker}(v)/\operatorname{Im}(u))\otimes_B(B/tB)
\end{equation}
\end{lemma}

Com efeito, a hipótese de que a sequência
\sref{IV.12.3.3.2} é exata implicará então
$(\operatorname{Ker}\,v/\operatorname{Im}\,u)\otimes_B(B/tB)=0$.
Como $t$ pertence ao ideal maximal de $B$ e
$\operatorname{Ker}(v)$ é um $B$-módulo de tipo finito
(pois $B$ é suposto noetheriano em \sref{IV.12.3.3.1}),
o lema de Nakayama provará que
$\operatorname{Ker}(v)=\operatorname{Im}(u)$.

Resta, pois, demonstrar \sref{IV.12.3.3.3}.
Ponhamos $u'=u\otimes1$, $v'=v\otimes1$,
$Z=\operatorname{Ker}(v)$ e $H=Z/\operatorname{Im}(u)$,
de modo que temos as sequências exatas
\begin{gather*}
0\to M\to Z\to H\to0\\
0\to Z\to N\xrightarrow{v}P\to0
\end{gather*}
Tensorizando por $B/tB$ e usando o lema
\sref{IV.3.4.1.4} e o fato de que $t$ é $P$-regular,
obtemos as sequências exatas
\begin{gather*}
M/tM\xrightarrow{w'}Z/tZ\to H/tH\to0\\
0\to Z/tZ\to N/tN\xrightarrow{v'}P/tP\to0
\end{gather*}
Daí, $\operatorname{Ker}(v')=Z/tZ$, e, como
$\operatorname{Im}(w')=\operatorname{Im}(u')$,
obtemos \sref{IV.12.3.3.4}.

\medskip
II) Suponhamos agora, além disso, que $\sh{L}_1$ seja
$f$-plano. Substituindo ainda, se necessário, $X$ por
$U$, podemos supor que $X=U$.
Trata-se de verificar que, para todo $x\in X$,
$(\sh{H}_1(\sh{L}_\bullet))_x=0$.
Ponhamos $s=f(x)$ e seja $\mathfrak{n}$ o ideal
$\mathfrak{m}_s\sh{O}_{X,x}$, que está contido no ideal
maximal de $\sh{O}_{X,x}$.
Como $(\sh{H}_1(\sh{L}_\bullet))_x$ é um
$\sh{O}_{X,x}$-módulo de tipo finito
(pois $X$ é localmente noetheriano), ele é separado
para a topologia $\mathfrak{n}$-ádica
\hyperref[0.7.3.5]{($\mathbf{0}_{\mathrm{I}}$,~7.3.5)}.
Basta, portanto, mostrar que seu \emph{completamento separado}
para essa topologia é $0$.
Ora, em virtude de \sref[III]{III.7.4.7.2}, esse
completamento separado é igual a
$\varprojlim_k\sh{H}_1(\sh{L}_\bullet\otimes_{\sh{O}_S}(\sh{O}_{S,s}/\mathfrak{m}_s^{k+1}))$.
Bastará, pois, provar que cada termo desse sistema
projetivo é nulo.
Isso é verdadeiro por hipótese para $k=0$;
raciocinemos, então, por indução sobre $k$.
A conclusão resultará do lema mais geral seguinte,
que aplicaremos com $A=\sh{O}_{S,s}/\mathfrak{m}_s^{k+1}$
e $\mathfrak{J}$ igual ao ideal maximal
$\mathfrak{m}_s/\mathfrak{m}_s^{k+1}$ de $A$.

\begin{lemma}[12.3.3.5]
\label{IV.12.3.3.5}
Sejam $A$ um anel, $\mathfrak{J}$ um ideal nilpotente de
$A$ e $L_\bullet:P\xrightarrow{u}Q\xrightarrow{v}R$
um complexo de $A$-módulos.
Ponhamos $A_0=A/\mathfrak{J}$ e
$L_\bullet^{(0)}=L_\bullet\otimes_A A_0:P_0\xrightarrow{u_0}Q_0\xrightarrow{v_0}R_0$.
Suponhamos que $\operatorname{gr}_{\mathfrak{J}}^{\bullet}(A)$
seja um $A_0$-módulo plano e que $Q$ e $R$ sejam
$A$-módulos planos.
Então, a relação $\mathrm{H}_1(L_\bullet^{(0)})=0$
implica $\mathrm{H}_1(L_\bullet)=0$.
\end{lemma}

Ponhamos
$L_\bullet^{(k)}=L_\bullet\otimes_A(A/\mathfrak{J}^{k+1}):P_k\xrightarrow{u_k}Q_k\xrightarrow{v_k}R_k$.
Como existe um inteiro $k\geqslant1$ tal que
$L_\bullet^{(k)}=L_\bullet$, bastará provar, por
indução sobre $k$, que a sequência
$P_k\xrightarrow{u_k}Q_k\xrightarrow{v_k}R_k$
\oldpage[IV-3]{187}
é exata; essa afirmação decorre da hipótese para $k=0$.
Seja, pois, $x\in Q_k$ um elemento tal que $v_k(x)=0$.
Como, pela hipótese de indução, a sequência
$P_{k-1}\xrightarrow{u_{k-1}}Q_{k-1}\xrightarrow{v_{k-1}}R_{k-1}$
é exata, a imagem canônica de $x$ em
$Q_{k-1}=Q_k/(\mathfrak{J}^kQ/\mathfrak{J}^{k+1}Q)$
pertence a $\operatorname{Im}(u_{k-1})$.
Existe, portanto, $z\in P_k$ tal que $x=u_k(z)+x'$,
com $x'\in\mathfrak{J}^kQ/\mathfrak{J}^{k+1}Q$.
Como $v_k\circ u_k=0$, a relação $v_k(x)=0$
implica $v_k(x')=0$.
Além disso, temos evidentemente
$v_k(x')\in\mathfrak{J}^kR/\mathfrak{J}^{k+1}R$.
Tudo se reduz, pois, a provar que a sequência
\[
\mathfrak{J}^kP/\mathfrak{J}^{k+1}P
\xrightarrow{u'}
\mathfrak{J}^kQ/\mathfrak{J}^{k+1}Q
\xrightarrow{v'}
\mathfrak{J}^kR/\mathfrak{J}^{k+1}R
\]
na qual $u'$ e $v'$ provêm de $u$ e $v$ por restrição
e passagem aos quocientes, é exata.
Ora, por hipótese, $\mathfrak{J}^k/\mathfrak{J}^{k+1}$
é um $(A/\mathfrak{J})$-módulo plano.
Logo, a sequência
\[
(P/\mathfrak{J}P)\otimes_{A/\mathfrak{J}}(\mathfrak{J}^k/\mathfrak{J}^{k+1})
\xrightarrow{u''}
(Q/\mathfrak{J}Q)\otimes_{A/\mathfrak{J}}(\mathfrak{J}^k/\mathfrak{J}^{k+1})
\xrightarrow{v''}
(R/\mathfrak{J}R)\otimes_{A/\mathfrak{J}}(\mathfrak{J}^k/\mathfrak{J}^{k+1})
\]
é exata.
Mas
$(Q/\mathfrak{J}Q)\otimes_{A/\mathfrak{J}}(\mathfrak{J}^k/\mathfrak{J}^{k+1})=Q\otimes_A(\mathfrak{J}^k/\mathfrak{J}^{k+1})$
identifica-se com $\mathfrak{J}^kQ/\mathfrak{J}^{k+1}Q$
em virtude da planicidade de $Q$ sobre $A$.
Do mesmo modo,
$(R/\mathfrak{J}R)\otimes_{A/\mathfrak{J}}(\mathfrak{J}^k/\mathfrak{J}^{k+1})$
identifica-se com $\mathfrak{J}^kR/\mathfrak{J}^{k+1}R$.
Finalmente, a imagem de $u'$ identifica-se com a
de $u''$, o que conclui a demonstração.
\end{proof}

\begin{corollary}[12.3.4]
\label{IV.12.3.4}
Sejam $f:X\to S$ um morfismo plano e localmente de apresentação
finita, e $\sh{F}$ e $\sh{G}$ dois $\sh{O}_X$-Módulos
de apresentação finita e $f$-planos.
Seja $n$ um inteiro $\geq0$ e seja $U$ o conjunto dos
$x\in X$ tais que
$(\mathscr{E}\!xt^n_{\sh{O}_{X_{f(x)}}}(\sh{F}_{f(x)},\sh{G}_{f(x)}))_x=0$
(respectivamente,
$(\mathscr{T}\!or_n^{\sh{O}_{X_{f(x)}}}(\sh{F}_{f(x)},\sh{G}_{f(x)}))_x=0$).
Então, $U$ é aberto e temos
$\mathscr{E}\!xt^n_{\sh{O}_X}(\sh{F},\sh{G})|_U=0$
(respectivamente, $\mathscr{T}\!or_n^{\sh{O}_X}(\sh{F},\sh{G})|_U=0$).
\end{corollary}

\begin{proof}
Podemos, evidentemente, restringir-nos ao caso em que $S$ e
$X$ são afins. Então, por \sref{IV.12.3.2}, existe uma
resolução à esquerda $\sh{L}_\bullet=(\sh{L}_i)$ de
$\sh{F}$ formada por $\sh{O}_X$-Módulos \emph{livres de tipo finito}.
Se pusermos $\sh{M}^\bullet=\mathscr{H}\!om(\sh{L}_\bullet,\sh{G})$
(respectivamente, $\sh{N}_\bullet=\sh{L}_\bullet\otimes_{\sh{O}_X}\sh{G}$),
deduziremos que cada $\sh{M}^i$ (respectivamente, $\sh{N}_i$)
é isomorfo a um $\sh{O}_X$-Módulo da forma $\sh{G}^m$.
Logo, os $\sh{M}^i$ e os $\sh{N}_i$ são $\sh{O}_X$-Módulos
de apresentação finita e $f$-planos.
Além disso, para toda mudança de base $S'\to S$, se pusermos
$X'=X\times_S S'$, $\sh{F}'=\sh{F}\otimes_S S'$,
$\sh{G}'=\sh{G}\otimes_S S'$ e
$\sh{L}'_\bullet=\sh{L}_\bullet\otimes_S S'$,
$\sh{L}'_\bullet$ ainda será uma resolução à esquerda de
$\sh{F}'$ por $\sh{O}_{X'}$-Módulos livres de tipo finito
\sref{IV.2.1.10}.
E $\sh{M}'^\bullet=\sh{M}^\bullet\otimes_S S'$ é igual a
$\mathscr{H}\!om(\sh{L}'_\bullet,\sh{G}')$
(respectivamente, $\sh{N}'_\bullet=\sh{N}_\bullet\otimes_S S'$
é igual a $\sh{L}'_\bullet\otimes_{\sh{O}_{X'}}\sh{G}'$),
pelo que precede.
Em particular, para todo $s\in S$, temos
$\mathscr{E}\!xt^n_{\sh{O}_{X_s}}(\sh{F}_s,\sh{G}_s)=\sh{H}^n((\sh{M}^\bullet)_s)$
(respectivamente,
$\mathscr{T}\!or_n^{\sh{O}_{X_s}}(\sh{F}_s,\sh{G}_s)=\sh{H}_n((\sh{N}_\bullet)_s)$).
Aplicando \sref{IV.12.3.3} aos complexos de Módulos
$f$-planos $\sh{M}^\bullet$ e $\sh{N}_\bullet$,
deduzimos imediatamente o corolário.
\end{proof}
